\documentclass[11pt,a4paper]{article}

% --- Jezik i kodiranje ---
\usepackage[T1]{fontenc}
\usepackage[utf8]{inputenc}
\usepackage[serbian]{babel}

% --- Layout i paketi ---
\usepackage[a4paper,top=1.5cm,bottom=1.8cm,left=1.8cm,right=1.8cm,headheight=28pt,headsep=10pt]{geometry}
\usepackage{xcolor}
\usepackage{listings}
\usepackage{fancyhdr}
\usepackage{lastpage}
\usepackage{enumitem}

% =========================================================
% Podaci o ispitu -- uredite po potrebi
% =========================================================
\newcommand{\univerzitet}{Univerzitet Crne Gore}
\newcommand{\fakultet}{Prirodno-matematički fakultet}
\newcommand{\studProgram}{Računarske nauke}
\newcommand{\predmet}{Strukture podataka}
\newcommand{\datumIspita}{20. maj 2026.}
\newcommand{\tipIspita}{Kolokvijum}
\newcommand{\trajanje}{80 minuta}
\newcommand{\ukupnoPoena}{20}

% =========================================================
% Stil za Java listinge
% =========================================================
\definecolor{javaKeyword}{RGB}{0,0,180}
\definecolor{javaString}{RGB}{0,128,0}
\definecolor{javaComment}{RGB}{120,120,120}
\definecolor{codeBg}{RGB}{248,248,248}
\definecolor{lineNum}{RGB}{140,140,140}

\lstdefinestyle{java}{
    language=Java,
    basicstyle=\ttfamily\small,
    keywordstyle=\color{javaKeyword}\bfseries,
    stringstyle=\color{javaString},
    commentstyle=\color{javaComment}\itshape,
    backgroundcolor=\color{codeBg},
    frame=single,
    rulecolor=\color{black!20},
    showstringspaces=false,
    tabsize=4,
    breaklines=true,
    captionpos=b,
    numbers=left,
    numberstyle=\tiny\color{lineNum},
    extendedchars=true,
    inputencoding=utf8,
    literate=%
        {č}{{\v{c}}}1 {ć}{{\'{c}}}1 {š}{{\v{s}}}1 {ž}{{\v{z}}}1 {đ}{{\dj}}1
        {Č}{{\v{C}}}1 {Ć}{{\'{C}}}1 {Š}{{\v{S}}}1 {Ž}{{\v{Z}}}1 {Đ}{{\DJ}}1,
}
\lstset{style=java}

% =========================================================
% Okruženje zadatak[poena] -- omogućava lstlisting unutra
% =========================================================
\newcounter{zadatak}
\newenvironment{zadatak}[1][]{%
    \par\vspace{0.9em}%
    \refstepcounter{zadatak}%
    \noindent\textbf{Zadatak \thezadatak.}%
    \def\tempPoena{#1}%
    \ifx\tempPoena\empty\else\ \textit{(#1 poena)}\fi%
    \par\medskip\noindent\ignorespaces%
}{\par}

% =========================================================
% Zaglavlje i podnožje
% =========================================================
\pagestyle{fancy}
\fancyhf{}
\fancyhead[L]{\small\univerzitet, \small\fakultet}
\fancyhead[R]{\small\predmet\ -- \tipIspita\\\small\datumIspita}
\fancyfoot[C]{\thepage\ / \pageref{LastPage}}
\renewcommand{\headrulewidth}{0.4pt}

% Stil za prvu stranicu -- isto zaglavlje kao na ostalim
\fancypagestyle{prva}{%
    \fancyhf{}%
    \fancyhead[L]{\small\univerzitet, \small\fakultet}%
    \fancyhead[R]{\studProgram}
    
    \fancyfoot[C]{\thepage\ / \pageref{LastPage}}%
    \fancyfoot[R]{\small\datumIspita}

    \renewcommand{\headrulewidth}{0.4pt}%
}

\begin{document}

\thispagestyle{prva}

\begin{center}
    {\Large\bfseries \predmet\  - \tipIspita}\\[0.2em]
    {\small Trajanje: \trajanje\ \quad$\bullet$\quad Ukupno: \ukupnoPoena\ poena}
\end{center}

\vspace{0.3em}
\noindent\textbf{Ime i prezime:}~\dotfill \qquad \textbf{Broj indeksa:}~\dotfill

\vspace{0.3em}
\hrule
\vspace{0.4em}

\noindent\textit{Uputstvo: Pažljivo pročitajte sve zadatke prije nego što počnete sa radom. Predajte izvorne fajlove pod imenima \texttt{Zadatak1.java}, \texttt{Zadatak2.java}, \ldots\ Dozvoljeno je korišćenje standardne Java biblioteke. Nije dozvoljeno korišćenje interneta niti komunikacija sa drugim kandidatima.}

\vspace{0.6em}

% =========================================================
% Zadaci -- uredite/dodajte po potrebi
% =========================================================

\begin{zadatak}[12]
Implementirati klasu \texttt{BST}, koja predstavlja binarno stablo traženja. U stablu se čuvaju cijeli brojevi. Potrebno je implementirati sljedeće metode:
\begin{enumerate}[label=(\alph*),leftmargin=2em,itemsep=0pt]
    \item (2 poena) dodavanje elementa u stablo,
    \item (2 poena) uklanjanje elementa iz stabla,
    \item (2 poena) štampanje čvorova stabla u inorder poretku,
    \item (4 poena) metodu \texttt{void nodeBalance()} koja za svaki čvor stabla štampa razliku dubine njegovog lijevog i desnog podstabla. Neophodno je voditi računa o optimalnosti algoritma koji nalazi disbalans. Za maksimalan broj bodova potrebno je implementirati algoritam složenosti $O(n)$,
    \item (2 poena) testirati napisane metode u \texttt{main} metodu.
\end{enumerate}
\end{zadatak}

\begin{zadatak}[8]
Implementirati klasu \texttt{DoubleStack}, koja predstavlja dva stacka pomoću jednog niza. Potrebno je omogućiti dodavanje elemenata u bilo koji od stackova sve dok se ne popuni niz. Implementirati sljedeće operacije:
\begin{enumerate}[label=(\alph*),leftmargin=2em,itemsep=0pt]
    \item (2 poena) \texttt{push(stack\_id, value)} -- dodaje element \texttt{value} na vrh stacka određenog parametrom \texttt{stack\_id},
    \item (2 poena) \texttt{pop(stack\_id)} -- uklanja i vraća element sa vrha stacka određenog parametrom \texttt{stack\_id},
    \item (2 poena) \texttt{peek(stack\_id)} -- vraća element sa vrha stacka određenog parametrom \texttt{stack\_id} bez njegovog uklanjanja,
    \item (2 poena) testirati napisane metode u \texttt{main} metodu.
\end{enumerate}
Sve metode treba da budu implementirane sa složenošću $O(1)$.
\end{zadatak}

\vfill
\noindent\rule{\linewidth}{0.4pt}\\[0.3em]
\textit{Ukupno: \ukupnoPoena\ poena.\quad Srećno!}

\end{document}
